\input{../tex-inputs/latex-leseansicht-vorspann}

\section[Felix und Ottilie Salten an Arthur Schnitzler, 12. 8. 1907]{L04250 Felix und Ottilie Salten an Arthur Schnitzler, 12. 8. 1907}
\nopagebreak\mylabel{L04250v}
\rehead{ }\normalsize\beginnumbering\briefempfaengerindex{Schnitzler, Arthur@\textsc{Schnitzler, Arthur}!zzzSalten, Ottilie@\emph{von Ottilie Salten}!1907-08-121@{12. 8. 1907}|(be}\briefempfaengerindex{Schnitzler, Arthur@\textsc{Schnitzler, Arthur}!zzzSalten, Felix@\emph{von Felix Salten}!1907-08-121@{12. 8. 1907}|(be}
\toendnotes[C]{\smallbreak\pagebreak[2]}
\correspDesc{Versand  durch Felix Salten, Ottilie Salten am 12. 8. 1907 in Marienbad
\newline{}Erhalt  durch Arthur Schnitzler im Zeitraum [13. 8. 1907 – 17. 8. 1907?] in Welsberg-Taisten}\toendnotes[C]{\smallbreak}
\Standort{CUL, Schnitzler, B 96a.}
\physDesc{Bildpostkarte, 222 Zeichen
\newline{}Handschrift: schwarze Tinte, lateinische Kurrent
\newline{}Versand: Stempel: »\nobreak{}\oindex{Marienbad@\textbf{Marienbad}|pwk}Marienbad, 12. VIII. 07, XI\nobreak{}«.  }\toendnotes[C]{\smallbreak}\pstart{}{\pb}Herrn D\textsuperscript{r} Arthur Schnitzler\pend{}\pstart{}Wildbad Waldbrunn\oindex{Wildbad Waldbrunn@\textbf{Wildbad Waldbrunn}, \emph{Spa}|pw}\pend{}\pstart{}bei Welsberg (Pustertal)
                  Tirol\oindex{Welsberg-Taisten@\textbf{Welsberg-Taisten}, \emph{Verwaltungsgebiet}|pw}\pend{}{\bigskip}
\pstart
           \textcolor{gray}{\textbf{{\pb}Forstwarte\oindex{Café Forstwarte@\textbf{Café Forstwarte}, \emph{Kaffeehaus}|pw}.}}\hfill \textcolor{gray}{\textbf{Marienbad\oindex{Marienbad@\textbf{Marienbad}|pw}.}}\pend
           \vspace{1em}
\pstart
           \noindent{}{\pb}Lieber, Danke herzlich für Ihren \label{K_L04250-1v}\edtext{Brief}{\lemma{\textnormal{\emph{Brief}}}\Cendnote{\textnormal{XXXX Auszeichnungsfehler: Dokument L03009 nicht gefunden.}}}\label{K_L04250-1}. Ich schreibe Ihnen morgen
               ausführlich. Inzwischen nur eilig viele Grüße von uns zu Ihnen.\pend
           \pstart Ihr \spacefill\mbox{Salten}\pend{}\selectlanguage{ngerman}\vspace{1em}
\pstart
           \noindent{}{[}hs. Salten:{]} Viele Grüße\pend
           \pstart \spacefill\mbox{Otti}\pend{}\selectlanguage{ngerman}\endnumbering\briefempfaengerindex{Schnitzler, Arthur@\textsc{Schnitzler, Arthur}!zzzSalten, Ottilie@\emph{von Ottilie Salten}!1907-08-121@{12. 8. 1907}|)be}\briefempfaengerindex{Schnitzler, Arthur@\textsc{Schnitzler, Arthur}!zzzSalten, Felix@\emph{von Felix Salten}!1907-08-121@{12. 8. 1907}|)be}\mylabel{L04250h}
\begin{anhang}
\end{anhang}\newcommand{\dateiname}{L04250}\newcommand{\titel}{Felix und Ottilie Salten an Arthur Schnitzler, 12. 8. 1907}\newcommand{\editorInnen}{Herausgegeben von Jahnke, SelmaMüller, Martin Anton}\input{../tex-inputs/latex-leseansicht-abspann}
